\subsection*{Additional Appendices}
\phantomsection
\addcontentsline{toc}{subsection}{Additional Appendices}

\subsubsection*{A.1\quad Wire Protocol Reference}
\phantomsection
\addcontentsline{toc}{subsubsection}{A.1\quad Wire Protocol Reference}
\label{sec:appendix-protocols}

This section provides the complete field specifications and representative implementation fragments for the three inter-component protocols described in Section~\ref{sec:wire-protocols}.

\paragraph{Complete llm-app to mcpd message reference.}

The \texttt{tool:exec} request carries the tool hash as a 64-character hexadecimal string.
If a prior \texttt{DEFER} decision returned a \texttt{ticket\_id}, the client includes it in a replay request:

\begin{lstlisting}[style=codestyle, language=json, caption={tool:exec request and response messages (llm-app $\to$ mcpd).}]
// tool:exec request
{ "kind": "tool:exec",  "req_id": 42,
  "session_id": "sess-abc", "agent_id": "agent-1",
  "app_id": "sys_app",  "tool_id": 1,
  "tool_hash": "a3f2...c91b",        // 64-char SHA-256 hex
  "payload": {"path": "/tmp"},
  "approval_ticket_id": null }       // set on replay after DEFER

// tool:exec response (ALLOW)
{ "req_id": 42, "status": "ok",
  "result": {"os": "Linux 6.8", "cpu": "4-core"},
  "decision": "ALLOW", "reason": "admitted",
  "ticket_id": 0,  "t_ms": 3 }

// tool:exec response (DEFER)
{ "req_id": 42, "status": "error",
  "result": {}, "error": "deferred pending approval",
  "decision": "DEFER", "reason": "risk_tag:sensitive_read",
  "ticket_id": 7,  "t_ms": 1 }
\end{lstlisting}

The \texttt{approval\_reply} message is issued by the client after the user confirms a deferred operation:

\begin{lstlisting}[style=codestyle, language=json, caption={approval\_reply message (llm-app $\to$ mcpd, issued after user confirms a deferred operation).}]
// approval_reply request
{ "sys": "approval_reply",  "session_id": "sess-abc",
  "ticket_id": 7,  "decision": "approve",
  "reason": "user confirmed",  "ttl_ms": 30000 }
\end{lstlisting}

\paragraph{Complete Netlink attribute encoding.}

All attributes are 4-byte aligned.
String attributes are NUL-terminated and passed as raw bytes.
The following Python fragment shows how \texttt{mcpd} encodes a \texttt{TOOL\_REQUEST} message:

\begin{lstlisting}[style=codestyle, language=Python, caption={Python fragment: encoding a \texttt{TOOL\_REQUEST} Generic Netlink message (mcpd $\to$ kernel).}]
import struct, socket, hashlib

FAMILY_NAME        = b"KERNEL_MCP\x00"
CMD_TOOL_REQUEST   = 10
ATTR_AGENT_ID,   ATTR_TOOL_ID     = 4,  2
ATTR_REQ_ID,     ATTR_TOOL_HASH   = 1,  19
ATTR_AGENT_BINDING, ATTR_AGENT_EPOCH = 28, 29

def build_nla(attr_type, data):
    # NLA header: 2-byte length, 2-byte type; 4-byte aligned
    hdr = struct.pack("HH", len(data) + 4, attr_type)
    pad = (4 - len(data) % 4) % 4
    return hdr + data + b"\x00" * pad

def encode_tool_request(agent_id, tool_id, req_id,
                        tool_hash, binding, epoch):
    attrs = (
        build_nla(ATTR_AGENT_ID,       agent_id.encode() + b"\x00") +
        build_nla(ATTR_TOOL_ID,        struct.pack("=I", tool_id))   +
        build_nla(ATTR_REQ_ID,         struct.pack("=Q", req_id))    +
        build_nla(ATTR_TOOL_HASH,      tool_hash.encode() + b"\x00") +
        build_nla(ATTR_AGENT_BINDING,  struct.pack("=Q", binding))   +
        build_nla(ATTR_AGENT_EPOCH,    struct.pack("=Q", epoch))
    )
    # Generic Netlink header: cmd (1 B), version (1 B), reserved (2 B)
    genl_hdr = struct.pack("BBH", CMD_TOOL_REQUEST, 1, 0)
    return genl_hdr + attrs
\end{lstlisting}

\paragraph{Semantic hash computation.}

The semantic hash used for tool identity is a full 256-bit SHA-256 digest computed over a canonical JSON serialisation of the following manifest fields: \texttt{tool\_id}, \texttt{name}, \texttt{app\_id}, \texttt{app\_name}, \texttt{risk\_tags}, \texttt{description}, \texttt{input\_schema}, \texttt{examples}, \texttt{path\_semantics}, and \texttt{approval\_policy}.
The \texttt{endpoint} field is deliberately excluded so that redeployment to a new socket path does not invalidate the kernel registration.

\begin{lstlisting}[style=codestyle, language=Python, caption={Python fragment: computing the manifest-derived semantic hash used for tool identity.}]
import json, hashlib

SEMANTIC_FIELDS = [
    "tool_id", "name", "app_id", "app_name",
    "risk_tags", "description", "input_schema",
    "examples", "path_semantics", "approval_policy",
]

def semantic_hash(tool_manifest: dict) -> str:
    canonical = {f: tool_manifest[f] for f in SEMANTIC_FIELDS}
    blob = json.dumps(canonical,
                      sort_keys=True,
                      separators=(",", ":"),
                      ensure_ascii=True).encode("utf-8")
    return hashlib.sha256(blob).hexdigest()  # 64-char hex; never truncated
\end{lstlisting}

The full 64-character digest is used without truncation.
A 32-bit prefix (8 hex characters) would be vulnerable to birthday collisions after approximately $2^{16}$ distinct tool registrations, which is a realistic size for a production deployment spanning multiple applications.

\paragraph{Complete manifest field specification.}

\begin{table}[htbp]
\centering
\footnotesize
\renewcommand{\arraystretch}{1.3}
\setlength{\tabcolsep}{6pt}
\begin{tabular}{l c p{7.8cm}}
\toprule
\textbf{Field} & \textbf{In hash?} & \textbf{Description} \\
\midrule
\multicolumn{3}{l}{\textit{Application-level fields}} \\[2pt]
\rowcolor{gray!8}
\texttt{app\_id}          & \ding{51} & Unique application identifier string \\
\texttt{app\_name}        & \ding{51} & Human-readable application name \\
\rowcolor{gray!8}
\texttt{transport}        & \ding{55} & Must be \texttt{"uds\_rpc"} \\
\texttt{endpoint}         & \ding{55} & Absolute socket path under \texttt{/tmp/linux-mcp-apps/} \\
\midrule
\multicolumn{3}{l}{\textit{Tool-level fields}} \\[2pt]
\rowcolor{gray!8}
\texttt{tool\_id}         & \ding{51} & Numeric identifier, unique within deployment \\
\texttt{name}             & \ding{51} & Tool name presented to the LLM planner \\
\rowcolor{gray!8}
\texttt{description}      & \ding{51} & Natural-language description of tool behaviour \\
\texttt{risk\_tags}       & \ding{51} & Risk labels controlling kernel DEFER (e.g., \texttt{"sensitive\_read"}) \\
\rowcolor{gray!8}
\texttt{input\_schema}    & \ding{51} & JSON Schema object for payload validation \\
\texttt{examples}         & \ding{51} & Example invocation payloads \\
\rowcolor{gray!8}
\texttt{path\_semantics}  & \ding{51} & Path interpretation hints (e.g., \texttt{"repo\_rel"}) \\
\texttt{approval\_policy} & \ding{51} & Policy name triggering kernel DEFER; null for unrestricted tools \\
\rowcolor{gray!8}
\texttt{operation}        & \ding{55} & Operation string forwarded to the tool service \\
\texttt{timeout\_ms}      & \ding{55} & Per-invocation execution timeout (operational metadata) \\
\bottomrule
\end{tabular}
\caption{Manifest field specification. \ding{51}~= included in the SHA-256 semantic identity hash; \ding{55}~= excluded (operational metadata whose change should not invalidate a tool registration). Fields are grouped by level: application-level fields appear once per \texttt{app\_id}; tool-level fields appear once per tool entry.}
\label{tab:manifest-fields}
\end{table}

\subsubsection*{A.2\quad System Comparison Table}
\phantomsection
\addcontentsline{toc}{subsubsection}{A.2\quad System Comparison Table}
\label{sec:appendix-comparison}

\begin{table}[htbp]
\centering
\footnotesize
\renewcommand{\arraystretch}{1.35}
\setlength{\tabcolsep}{5pt}
\begin{tabular}{p{2.0cm} p{2.8cm} p{2.6cm} p{2.6cm} p{4.2cm}}
\toprule
\textbf{System}
  & \textbf{Control-plane placement}
  & \textbf{Tool identity}
  & \textbf{State durability}
  & \textbf{Trust anchor \& security boundary} \\
\midrule

\rowcolor{gray!12}
\textbf{\texttt{linux-mcp}} \newline\textbf{(this work)}
  & Kernel module + user-space gateway
  & SHA-256 of JSON manifest; kernel-verified at every invocation
  & Approval tickets and registry survive daemon crash (kernel-held)
  & \textbf{Kernel}; enforcement is independent of gateway correctness \\

\texttt{OpenClaw} \cite{openclaw2026}
  & User-space gateway (persistent process)
  & Skill name/version string in gateway registry
  & Lost on process restart
  & User-space process; bypassable via logic flaws \cite{openclaw_security1, openclaw_security2} \\

\rowcolor{gray!12}
Harness Agent \cite{harness2025}
  & Cloud-hosted pipeline orchestrator
  & Pipeline step + RBAC connector identity
  & Cloud-hosted; platform-resilient
  & SaaS control plane; RBAC-based governance \\

AIOS / AgentOS \cite{mei2024aios, liu2026agentos}
  & New OS-level agent kernel (proposed)
  & Agent/tool registration in LLM scheduler
  & Not evaluated (proposal stage)
  & New OS substrate; full redesign \\

\rowcolor{gray!12}
LangChain / LangGraph \cite{langchain2022, langgraph2024}
  & In-process user-space library
  & Tool name string; no cryptographic binding
  & Lost on process crash
  & Application process; no inter-invocation isolation \\

Apple Intelligence \cite{apple2024pcc}
  & On-device OS + Private Cloud Compute
  & OS capability entitlements
  & High; OS-persistent services
  & Hardware attestation + OS entitlements; closed ecosystem \\

\bottomrule
\end{tabular}
\caption{Comparison of representative agent systems. The key differentiator for \texttt{linux-mcp} is the combination of cryptographic tool identity and kernel-held durable state, both absent from all user-space frameworks.}
\label{tab:system_comparison}
\end{table}
